\documentclass[11pt]{article}
\usepackage[letterpaper,margin=0.9in]{geometry}
\usepackage{amsmath,amssymb,amsthm}
\usepackage{microtype}
\usepackage[colorlinks=true,linkcolor=blue,urlcolor=blue]{hyperref}
\newtheorem{proposition}{Proposition}
\newcommand{\HH}{\mathcal H}
\newcommand{\LL}{\mathcal L}
\setlength{\parindent}{0pt}
\setlength{\parskip}{7pt}
\begin{document}
\begin{center}
{\Large\bfseries Theorem 1.1: formal statement\\[3pt] and initial verified lemmas}\\[9pt]
Falconer distance-gauge formalization project\\
October 8, 2026
\end{center}

\textbf{Status.} This is an incomplete formalization notebook. The main theorem
is an explicit proof obligation; it has not passed Comparator and is not
registered on Palomar. The results below record the proved foundation.

\section*{The exact target}
The uploaded manuscript defines
\[
h_\theta(r)=r\exp\!\left(-\bigl(\log(1/r)\bigr)^\theta\right)
\quad(0<r<1),\qquad h_\theta(0)=0.
\]
Its Theorem 1.1 asserts that, for a compact set $E\subseteq\mathbb R^2$,
\[
\frac23<\theta<1,\qquad \HH^{h_\theta}(E)>0
\quad\Longrightarrow\quad \LL^1(\Delta(E))>0.
\]
\[
\Delta(E)=\{\lVert x-y\rVert:x,y\in E\}.
\]
The formalization uses the Euclidean norm and Mathlib's arbitrary-gauge
Hausdorff measure constructor. For radii at least one, the gauge is extended
by the identity; this does not affect its small-radius Hausdorff measure.
The covering characterization has been proved:
\[
\HH^{h_\theta}(E)=
\sup_{\rho>0}\ \inf_{\substack{E\subseteq\bigcup_{n\ge0}U_n\\
\operatorname{diam}(U_n)\le\rho}}
\sum_{\substack{n\ge0\\U_n\ne\varnothing}}
h_\theta\bigl(\operatorname{diam}(U_n)\bigr).
\]
The independent \texttt{Challenge.lean} imports only Mathlib and repeats all
custom definitions. The solution does not import Challenge. The theorem
statement has no auxiliary Frostman, projection, or energy hypothesis.

\section*{The geometric foundation}
\begin{proposition}
For $0<r<1$, the gauge is positive and satisfies $h_\theta(r)\le r$.
Consequently $\HH^{h_\theta}(E)\le\HH^1(E)$, and countable sets have zero
gauge measure. Positive gauge measure supplies two distinct points of $E$.
\end{proposition}
\begin{proof}
The exponential is strictly positive. The logarithm and its real power are
nonnegative, so the exponential is at most one. Mathlib's gauge comparison
theorem gives the measure inequality. Singletons are null for $\HH^1$,
hence also for the gauge measure; countable subadditivity gives the countable
case. A set containing at most one point is countable.
\end{proof}
The all-pairs distance map is continuous. It therefore maps the compact
product $E\times E$ to a compact, Lebesgue-measurable distance set. Two distinct
points give a positive distance; that fact alone does not give positive length.

\newpage
\section*{The scalar summability foundation}
\begin{proposition}
For every $s,c>0$, one has $\log n\le c n^s$ for all sufficiently large
natural numbers $n$. For every $\theta,c>0$,
\[
\sum_{n=0}^{\infty}\exp(-c n^\theta)<\infty.
\]
\end{proposition}
\begin{proof}
Mathlib proves $\log x=o(x^s)$ as $x\to\infty$ for $s>0$. Restricting its
eventual norm bound to natural numbers proves the first assertion. Applying
that assertion with coefficient $c/2$ gives
\[
\log n\le\frac c2 n^\theta
\quad\Longrightarrow\quad
\exp(-c n^\theta)\le\exp(-2\log n)=n^{-2}.
\]
The inverse-square series converges. The omitted finite initial segment has
no effect on summability, so comparison proves the claim.
\end{proof}

\begin{proposition}
The branch-counting exponent and filter-error power obey
\[
3\theta-2>0\ \Longleftrightarrow\ \theta>\frac23,
\qquad
\theta>\frac23\ \Longrightarrow\ -8(2\theta-1)<-1.
\]
Consequently, in the theorem's parameter range, every positive coefficient
eventually absorbs the logarithm into $n^{3\theta-2}$, and
\[
\sum_{n=0}^{\infty}n^{-8(2\theta-1)}<\infty.
\]
\end{proposition}
\begin{proof}
The exponent inequalities are linear arithmetic. Apply the logarithmic
domination result with exponent $3\theta-2$, and apply Mathlib's real-power
series criterion to the filter-error exponent.
\end{proof}

These results justify scalar convergence once suitable shell energies and
filter errors are supplied. They do not establish the geometric estimates
that produce those bounds.

\section*{Remaining analytic work and verification}
The proof still needs the gauge Frostman construction, logarithmically
weighted Fourier and projection estimates, endpoint Orlicz radial projection
control, the summable Fourier reconstruction criterion, quantitative regular
decomposition, the four uniform energy estimates, and the multiscale budget
induction. Fixed-parameter estimates require new uniform bounds before they
can be used with parameters that shrink as the terminal scale grows.

The main theorem currently has one explicit proof hole. Comparator permits
the three standard axioms and must reject the target's proof hole.
A successful ordinary build of this skeleton is not a successful theorem
verification.

The public project includes a separate sandboxed Comparator job and a manual
full Palomar preflight. Registration requires a completed proof, verification
of the exact public commit, truthful metadata, and the registry review.

\smallskip
\href{https://github.com/CoolRmal/falconer-theta-gauge}{Public project: CoolRmal/falconer-theta-gauge}
\end{document}
